\section{Radio Line Stacking}
\label{sec:line_stacking}

% (for a deeper look at line strength based on GDIGGS see \url{https://www.overleaf.com/project/68b0918d19fe364fd7b8db66})
The principle of line stacking is that each RRL is a measurement of the same underlying physical quantities, namely the emission measure and electron temperature of the ionized gas. We model each observation \(x_i\) as an independent draw from a normal distribution, \(x_i \sim \mathcal{N}(T_L, \sigma_i^2)\), where \(T_L\) is the line strength and \(\sigma_i\) is the brightness temperature sensitivity\footnote{Brightness temperature sensitivity is equivalent to the noise in the measurement and the two terms are used interchangeably} of line \(i\). In the ideal case, for \(n\) independent samples with mean \(\mu\) and common variance \(\sigma^2\), the variance on the mean is \(\sigma^2/n\). For RRLs, however, lines at different frequencies have different noise, so the optimal estimator is not a simple average but a weighted average. \\

The inverse-variance-weighted average minimizes the variance of the mean when variances are not constant. The weighted mean \(\hat{\mu}\) and variance \(\mathrm{Var}(\hat{\mu})\) are given by
\begin{equation}
    \hat{\mu} = \frac{\sum_i x_i w_i}{\sum_i w_i}, \quad \mathrm{Var}(\hat{\mu}) = \frac{\sum_i w_i^2 \sigma_i^2}{(\sum_i w_i)^2}
    \label{eq:weighted_avg}
\end{equation}
\begin{equation}
    w_i = \frac{1}{\sigma_i^2}.
\end{equation}
When all variances are equal, equation \eqref{eq:weighted_avg} reduces to the simple arithmetic mean.

Inverse-variance-weighted average assumes that each sample has the same mean, however RRLs have different expected line strengths at different frequencies, meaning the expected mean \(T_L\) is not constant. To account for this we rescale each line and its noise to a common reference frequency \(\nu_0\). After rescaling, each observation has the same expected mean but different variances. The rescaled line strength \(x'\), noise \(\sigma'\) and weights \(w'\) are given by

\begin{equation}
    x_i' = x_i \frac{\nu_i}{\nu_0}, \quad \sigma_i' = \sigma_i \frac{\nu_i}{\nu_0}, \quad w_i' = \frac{\nu_0^2}{\sigma_i^2 \nu_i^2}.
    \label{eq:rescaled_signal_and_noise}
\end{equation}
% The inverse-variance-weighted mean, variance and standard deviation of the rescaled quantities are given by
% \begin{equation}
%     \hat{\mu}'
%     = \frac{\sum_i x_i' w_i'}{\sum_i w_i'}
%     % = \frac{\sum_i (x_i\frac{\nu_i}{\nu_0}) ( \frac{\nu_0^2}{\sigma_i^2\nu_i^2} )} {\sum_i \frac{\nu_0^2}{\sigma_i^2\nu_i^2}}
%     = \frac{1}{\nu_0} \frac{\sum_i x_i \frac{1}{\sigma_i^2\nu_i}}{\sum_i \frac{1}{\sigma_i^2\nu_i^2}}
%     \label{eq:stacked_weighted_average}
% \end{equation}
% \begin{equation}
%     \mathrm{Var}(\hat{\mu}')
%     = \frac{\sum_i w_i'^2 \sigma_i'^2}{(\sum_i w_i')^2}
%     % = \frac{\sum_i (\frac{\nu_0^2}{\sigma_i^2 \nu_i^2})^2 (\sigma_i \frac{\nu_i}{\nu_0})^2}{(\sum_i \frac{\nu_0^2}{\sigma_i^2 \nu_i^2} )^2}
%     % = \frac{1}{\nu_0^2} \frac{\sum_i \frac{1}{\sigma_i^2\nu_i^2}}{(\sum_i \frac{1}{\sigma_i^2 \nu_i^2})^2}
%     = \left( \nu_0^2 \sum_i \frac{1}{\sigma_i^2 \nu_i^2} \right)^{-1}
%     \label{eq:stacked_variance}
% \end{equation}
% \begin{equation}
%     \hat{\sigma}' = \sqrt{\mathrm{Var}(\hat{\mu}')} = \nu_0^{-1} \left( \sum_i \frac{1}{\sigma_i^2 \nu_i^2} \right)^{-1/2}.
%     \label{eq:stacked_std}
% \end{equation}
The choice of \(\nu_0\) is arbitrary and adds a scaling factor to the estimated mean. The SNR is independent of the choice of \(\nu_0\) and is given by
\begin{equation}
    \mathrm{SNR} = \frac{\hat{\mu}'}{\sqrt{\mathrm{Var}(\hat{\mu}')}} = \frac{\sum_i x_i \frac{1}{\sigma_i^2 \nu_i}}{\sqrt{\sum_i \frac{1}{\sigma_i^2 \nu_i^2}}}.
\end{equation}
